\subsection{Adaptive Inverse Kinematics}\label{subsec:aik}

\subsubsection{Forward vs.\ Inverse Kinematics}

\emph{Forward kinematics} (\acrshort{fk}) maps a joint-configuration vector
to joint positions: given $\boldsymbol{\theta}$, every joint's world position
is obtained by composing rigid-body transforms along the kinematic chain.
\emph{Inverse kinematics} (\acrshort{ik}) poses the dual problem: given target
positions for a subset of joints, recover the configuration that produces them.
Classical \acrshort{ik} solvers --- FABRIK, cyclic-coordinate descent,
Jacobian-based methods --- iterate until convergence, which makes them
sensitive to initialisation and frame rate.

When every joint position is observed simultaneously, the system is
\emph{fully observed} and the \acrshort{ik} problem can be solved in closed
form without iteration. The Adaptive \acrshort{ik} algorithm
(AIK) exploits this property for the hand skeleton:
the global root orientation is recovered via a single singular value
decomposition (\acrshort{svd}), and every subsequent joint rotation follows
analytically from a swing--twist decomposition along the kinematic chain.

\subsubsection{Global Root Orientation via Point-Set Alignment}

The global orientation of the root joint is the rotation $R_0 \in SO(3)$ that
best aligns two corresponding sets of vectors --- those derived from the
template skeleton and those derived from the observed targets.
Let $\mathbf{A}, \mathbf{B} \in \mathbb{R}^{3 \times N}$ be column-stacked
matrices of $N$ corresponding vectors in the template and target,
respectively. The optimal rotation minimises the sum of squared residuals:
\begin{equation}
  R_0 \;=\; \arg\min_{R \in SO(3)}
  \sum_{i=1}^{N} \bigl\|\mathbf{b}_i - R\,\mathbf{a}_i\bigr\|^2.
\end{equation}
Following Arun et al.\ \cite{arun1987leastsquares}, the solution is obtained
from the cross-covariance matrix
\begin{equation}\label{eq:H}
  \mathbf{H} \;=\; \mathbf{A}\mathbf{B}^{\top} \;\in\; \mathbb{R}^{3 \times 3},
\end{equation}
whose \acrshort{svd} $\mathbf{H} = \mathbf{U}\mathbf{S}\mathbf{V}^{\top}$
yields
\begin{equation}\label{eq:R0}
  R_0 \;=\; \mathbf{V}\mathbf{U}^{\top}.
\end{equation}
A proper rotation requires $\det R_0 = +1$. When the input vectors are
coplanar, the smallest singular value of $\mathbf{S}$ vanishes and
$\det(\mathbf{V}\mathbf{U}^{\top})$ may equal $-1$ (a reflection). In that
degenerate case, the sign of the last column of $\mathbf{V}$ is negated before
recomputing $R_0$, recovering the nearest element of $SO(3)$, as prescribed by
Arun et al.\ \cite{arun1987leastsquares}.

For the hand skeleton, $\mathbf{A}$ and $\mathbf{B}$ are formed from the five
vectors pointing from the wrist to each MCP joint in the
template and in the target, respectively ($N = 5$). Because the metacarpal
plate is anatomically near-rigid, all five MCP joints are assigned the same
global rotation $R_0$.

\subsubsection{Per-Joint Rotation via Kinematic Chain Traversal}

For each remaining joint $k$, processed in parent-before-child order through
the kinematic tree, the local rotation $R_{p,k}$ relative to its parent $p =
\mathrm{pa}(k)$ is derived as follows. Let $\mathbf{T}_j$ denote the $j$-th
joint position in the template (rest pose), $\mathbf{q}_j$ the corresponding
computed world position, and $\mathbf{P}_k$ the observed target position of
joint $k$.

\textbf{Step~1 --- FK propagation.}
The parent's world position is updated by applying its global rotation to the
template bone and translating from the grandparent $p' = \mathrm{pa}(p)$:
\begin{equation}
  \mathbf{q}_{p} \;=\;
  R_{p'}\,\bigl(\mathbf{T}_{p} - \mathbf{T}_{p'}\bigr) + \mathbf{q}_{p'}.
\end{equation}

\textbf{Step~2 --- local-frame projection.}
The target is brought into the parent's local coordinate frame:
\begin{equation}
  \delta\mathbf{p}_k \;=\; R_{p}^{-1}
  \bigl(\mathbf{P}_k - \mathbf{q}_{p}\bigr).
\end{equation}

\textbf{Step~3 --- template bone vector.}
The rest-pose direction of the bone is
\begin{equation}
  \delta\mathbf{t}_k \;=\; \mathbf{T}_k - \mathbf{T}_{p}.
\end{equation}

\textbf{Step~4 --- swing axis and angle.}
The swing rotation is the unique rotation in the plane spanned by
$\delta\mathbf{t}_k$ and $\delta\mathbf{p}_k$ that maps the former onto the
latter:
\begin{equation}\label{eq:swing}
  \hat{\mathbf{n}} \;=\;
  \frac{\delta\mathbf{t}_k \times \delta\mathbf{p}_k}
       {\|\delta\mathbf{t}_k \times \delta\mathbf{p}_k\|},
  \qquad
  \alpha \;=\; \arccos\!\left(
    \frac{\delta\mathbf{t}_k \cdot \delta\mathbf{p}_k}
         {\|\delta\mathbf{t}_k\|\,\|\delta\mathbf{p}_k\|}
  \right).
\end{equation}

\textbf{Step~5 --- swing--twist decomposition.}
Any 3D rotation can be decomposed into a \emph{swing} component (rotation
perpendicular to a reference axis) and a \emph{twist} component (rotation
about that axis). Taking the bone direction $\delta\mathbf{t}_k$ as the twist
axis, the local rotation is
\begin{equation}\label{eq:local_rot}
  R_{p,k} \;=\;
  D_{\mathrm{sw}}(\hat{\mathbf{n}},\,\alpha)
  \cdot
  D_{\mathrm{tw}}(\delta\mathbf{t}_k,\,\varphi_k),
  \qquad
  R_k \;=\; R_{p} \cdot R_{p,k},
\end{equation}
where $D_{\mathrm{sw}}$ and $D_{\mathrm{tw}}$ are axis-angle rotation matrices.
The AIK algorithm sets $\varphi_k = 0$ for all joints, so
$D_{\mathrm{tw}} = \mathbf{I}$ and the local rotation reduces to the swing
alone. This is the central theoretical simplification: axial spin about a
finger bone is physiologically negligible, and omitting it removes the
directional ambiguity of the cross product when the template and target bones
are nearly collinear.

The traversal yields, for every joint in the tree, a local rotation matrix
$R_{p,k}$ relative to its parent and an absolute rotation matrix $R_k$ in
world space, together constituting a complete pose of the skeleton.
